\documentclass[../main.tex]{subfiles}

\graphicspath{{\subfix{../images/}}}

\newpage

\begin{document}

\section{Filtri attivi}

\subsection{Filtri del I ordine}
\subsubsection{Integratore}
\subsubsection{Passa-basso}
\subsubsection{Derivatore e passa-alto}
\subsubsection{Passa-banda}
\subsubsection{Rotatore di fase}

\subsection{Filtri del II ordine}
\subsubsection{Passa-basso}
\subsubsection{Passa-alto}
\subsubsection{Passa-banda}
\subsubsection{Elimina-banda}
\subsubsection{Passa tutto o giratore}

\subsection{Circuiti per filtri del II ordine}
\subsubsection{Celle a guadagno finito}
\subsubsection{Celle a guadagno infinito}
\subsubsection{Celle basate su amplificatori operazionali multipli}

\subsection{Filtri di ordine superiore al II}
\subsubsection{Maschera di progetto }
\subsubsection{Risposte standard}
\subsubsection{Progetto di un filtro passa basso}
\subsubsection{Circuito di simulazione di un’induttanza}
\subsubsection{Dati per il progetto di filtri passa-basso}

\subsection{Filtri a capacità commutate}
\subsubsection{Principio di funzionamento}
\subsubsection{Principio di funzionamento}
\subsubsection{Integratore}
\subsubsection{Limiti di frequenza di clock}
\subsubsection{Effetti delle capacità parassite}
\subsubsection{Integratori stray insensitive}
\subsubsection{Comportamento in frequenza}
\subsubsection{Filtro del secondo ordine con cella biquadratica}
\subsubsection{Approfondimenti}

\end{document}